\begin{frame}{Recap: Where 3D Has Already Appeared}
  \begin{itemize}\itemsep0.3em
    \item \textbf{Covered earlier:} Depth Anything predicts a depth map from a single image, and SuperGlue matches keypoints between two images. Both were used as pretrained components. This lecture supplies the geometry that turns such outputs into cameras, depths, and maps.
    \item \textbf{Also covered earlier:} optical flow as a dense motion field, RGB-D cameras and LiDAR as depth sensors, and image retrieval with global descriptors. Each returns here as one stage of a 3D pipeline.
    \item \textbf{Backbones we reuse:} ViT and cross-attention, DINOv2 features, masked autoencoders, and Stable Diffusion. The learned depth and 3D models later in the lecture are assembled from these parts.
    \item \textbf{Route:} camera model $\rightarrow$ two views $\rightarrow$ stereo $\rightarrow$ single-image depth $\rightarrow$ pose and structure from motion $\rightarrow$ SLAM $\rightarrow$ feed-forward 3D models.
  \end{itemize}
\end{frame}
